\section{The spectral gap at a large observation time}
\label{sec:terminal}

We complete the observation argument by proving a gap after a sufficiently
large fixed observation time. The deterministic two-spin curvature estimate
is uniform over external fields. The final probabilistic statement instead
concerns the specified observation law, at fixed $j<1$.
We adapt the signed two-spin method of Wang~\cite[Lemma~3.4]{Wang2026},
retaining the signed first-order interaction and estimating the remaining
terms by an entrywise-square matrix and a quartic row error. The complete
calculation below proves the precise variant used here.

For a symmetric matrix $J$ with zero diagonal, define the symmetric matrices
\[
 U_{ij}=\tanh J_{ij},\qquad Q_{ij}=U_{ij}^2,\qquad
 r_4(U)=\max_i\sum_j U_{ij}^4.
\]
In particular, $U_{ii}=Q_{ii}=0$.  The square defining $Q$ is entrywise.
At a field $h_0$, put
\[
 Hf=\sum_i(I-P_i)f
    =\sum_i(x_i-m_i^1)d_i f,\qquad
 g_i=d_i f,\qquad a_i=v_i^1g_i.
\]
Each $P_i$ is an orthogonal projection in $L^2(\mu_{h_0})$, and therefore
$H$ is self-adjoint and positive, with
\begin{equation}
 \E_{\mu_{h_0}}[fHf]=\D_{h_0}(f).
 \label{eq:terminal-generator-form}
\end{equation}

\begin{lemma}[Signed curvature inequality]
\label{lem:curvature}
There are universal constants $u_0>0$ and $C<\infty$ such that, for every
$n$, every symmetric zero-diagonal $J$ satisfying
$\max_{i,j}|U_{ij}|\le u_0$, every $h_0\in\R^n$, and every real function $f$
on the discrete cube,
\begin{equation}
 \E_{\mu_{h_0}}[(Hf)^2]
 \ge (1-Cr_4(U))\D_{h_0}(f)
       -\E_{\mu_{h_0}}\bigl[a^{\mathsf T}Ua
                         +C|a|^{\mathsf T}Q|a|\bigr].
 \label{eq:terminal-curvature}
\end{equation}
Here $|a|$ denotes coordinatewise absolute value.
\end{lemma}

\begin{proof}
The curvature bound is obtained by summing an exact two-spin calculation.
Fix $i\ne j$, condition on the other spins, and call the remaining spins
$x,y$. Their conditional density is proportional to
$\exp(\alpha x+\gamma y+J_{ij}xy)$. Introduce the scalar parameters
\[
 a=\tanh\alpha,\quad b=\tanh\gamma,\quad
 A=1-a^2,\quad B=1-b^2,\quad u=\tanh J_{ij}.
\]
The letters $a,b,A,B$ have these scalar meanings only during the pair
calculation. Use $\E_0$ for the product law of means $a,b$, and
$\E_{ij}$ for the interacting conditional law. The identity
$e^{J_{ij}xy}=\cosh(J_{ij})(1+uxy)$ gives their exact density ratio:
\begin{equation}
 \frac{\dd\mu_{ij}}{\dd\mu_{0,ij}}(x,y)
       =\frac{1+uxy}{1+uab}.
 \label{eq:pair-density}
\end{equation}
Expand $f$ in the four-dimensional space of functions of $x,y$. Its
two gradients then have the unique representation
\[
 g_i=p+r(y-b),\qquad g_j=s+r(x-a)
\]
with real $p,s,r$. The shared coefficient $r$ is the mixed-spin term
of this expansion. The conditional means and variances satisfy
\begin{align}
 x-m_i^1&=\frac{(x-a)(1-uxy)}{1+uay},&
 v_i^1&=\frac{A(1-u^2)}{(1+uay)^2},\notag\\
 y-m_j^1&=\frac{(y-b)(1-uxy)}{1+ubx},&
 v_j^1&=\frac{B(1-u^2)}{(1+ubx)^2}.
 \label{eq:pair-conditional-identities}
\end{align}

For a function $F$ of one spin, write $F(z)=\overline F+z\,dF$, and set
\[
 F_1(y)=\frac{g_i(y)}{1+uay},\qquad
 F_2(x)=\frac{g_j(x)}{1+ubx}.
\]
The identity $(1-uxy)^2(1+uxy)=(1-u^2)(1-uxy)$ and independence under
$\E_0$ give the following exact expression for the cross term:
\begin{equation}
 \frac{\E_{ij}[(x-m_i^1)g_i\,(y-m_j^1)g_j]}{AB}
 =\frac{1-u^2}{1+uab}
       \bigl(dF_1dF_2-u\overline F_1\overline F_2\bigr).
 \label{eq:pair-cross-exact}
\end{equation}
Both terms follow from the one-spin identities
$\E_0[(x-a)F_2(x)]=A\,dF_2$ and
$\E_0[(x-a)xF_2(x)]=A\overline F_2$, together with their $y$
counterparts. Thus the cross term has been reduced to four scalar
coefficients.

The four coefficients appearing here are
\begin{align*}
 dF_1&=\frac{r(1+uab)-uap}{1-u^2a^2},&
 \overline F_1&=\frac{p-(b+ua)r}{1-u^2a^2},\\
 dF_2&=\frac{r(1+uab)-ubs}{1-u^2b^2},&
 \overline F_2&=\frac{s-(a+ub)r}{1-u^2b^2}.
\end{align*}
Consequently the right-hand side of \eqref{eq:pair-cross-exact} is
\begin{equation}
 K\bigl[(-u+u^2ab)ps+u^2bA\,rp+u^2aB\,rs+Rr^2\bigr],
 \label{eq:pair-cross-coefficients}
\end{equation}
where
\begin{align*}
 K&=\frac{1-u^2}{(1+uab)(1-u^2a^2)(1-u^2b^2)},\\
 R&=(1+uab)^2-u(b+ua)(a+ub)\\
  &=1+uab+u^2(a^2b^2-a^2-b^2)-u^3ab.
\end{align*}
The mixed coefficients $rp$ and $rs$ begin at order $u^2$: their
first-order terms have canceled. Keeping this cancellation is what will
leave only a quartic row cost after Young's inequality.

For $|u|\le1/4$ and $|a|,|b|\le1$, all denominators above are bounded
away from zero.  There is thus a universal constant $C_1$ such that
\[
 |K-1|+|R-1|\le C_1|u|.
\]
Consequently the $ps$ coefficient in \eqref{eq:pair-cross-coefficients}
is $-u+O(u^2)$, both mixed coefficients are $O(u^2)$, and the $r^2$
coefficient is $1+O(|u|)$. These constants are uniform in the conditional
means $a,b$: all the displayed rational functions and their required
derivatives are bounded on the compact parameter range.

Return temporarily to the notation $a_i=v_i^1g_i$ and $a_j=v_j^1g_j$ for
the two gradient weights.  Equations \eqref{eq:pair-density} and
\eqref{eq:pair-conditional-identities} imply
\begin{equation}
 \frac{\E_{ij}[a_i a_j]}{AB}
       =\E_0\bigl[w_u(x,y)g_i(y)g_j(x)\bigr],\qquad
 w_u(x,y)=\frac{(1-u^2)^2(1+uxy)}
 {(1+uab)(1+uay)^2(1+ubx)^2}.
 \label{eq:pair-gradient-product}
\end{equation}
Uniformly in $a,b,x,y$, one has
\[
 c_1\le w_u(x,y)\le C_2,\qquad
 |w_u(x,y)-1|\le C_2|u|\qquad (|u|\le1/4)
\]
for universal positive constants $c_1,C_2$.  Expanding
$g_i(y)g_j(x)$ and using that the centered factors have zero product-law
mean gives
\begin{equation}
 \frac{\E_{ij}[a_i a_j]}{AB}
       =ps+e_0ps+e_1rp+e_2rs+e_3r^2,
       \qquad |e_\ell|\le C_3|u|.
 \label{eq:pair-gradient-coefficients}
\end{equation}
To obtain the four error bounds, expand the product into its $ps$,
$rp$, $rs$, and $r^2$ terms. Each centered spin is bounded by two, and
the uniform estimate for $w_u-1$ controls every resulting coefficient.

Adding $u$ times \eqref{eq:pair-gradient-coefficients} to
\eqref{eq:pair-cross-coefficients}, and decreasing the universal $u_0>0$
if necessary, yields
\begin{align*}
 &\frac{\E_{ij}[(x-m_i^1)g_i\,(y-m_j^1)g_j]
                    +u\E_{ij}[a_i a_j]}{AB}\\
 &\hspace{15mm}\ge
       \frac12r^2-C_4u^2|ps|
                       -C_4u^2|r|(|p|+|s|)\\
 &\hspace{15mm}\ge -C_5u^2|ps|-C_5u^4(p^2+s^2).
\end{align*}
The last inequality is Young's inequality, absorbing the mixed terms
into $r^2/2$.

We must now express the two scalar costs in the heat-bath weights that
can be summed over pairs. The lower bound for $w_u$ in
\eqref{eq:pair-gradient-product}, followed by product independence,
gives
\begin{align*}
 \E_{ij}|a_i a_j|
 &\ge c_1AB\,\E_0|g_i(y)g_j(x)|\\
 &=c_1AB\,\E_0|g_i(y)|\,\E_0|g_j(x)|
 \ge c_1AB|ps|.
\end{align*}
Similarly, the density in \eqref{eq:pair-density} and the ratios
$v_i^1/A$, $v_j^1/B$ are bounded above and below by universal positive
constants for $|u|\le u_0$.  Hence
\[
 \E_{ij}[v_i^1g_i^2]\ge c_2A\E_0g_i^2
                       \ge c_2Ap^2\ge c_2ABp^2,
 \qquad
 \E_{ij}[v_j^1g_j^2]\ge c_2ABs^2.
\]
The comparison constants do not deteriorate when a conditional variance
$A$ or $B$ is small. Substituting these bounds into the preceding scalar
inequality proves the pair estimate
\begin{align}
 \E_{ij}[(x-m_i^1)g_i\,(y-m_j^1)g_j]
 \ge{}&-u\E_{ij}[a_i a_j]-C_6u^2\E_{ij}|a_i a_j|\notag\\
 &-C_6u^4\E_{ij}[v_i^1g_i^2+v_j^1g_j^2].
 \label{eq:terminal-pair-bound}
\end{align}

To assemble the curvature form, first condition on all but spin $i$ in
the diagonal term of $\E(Hf)^2$. Their sum is
$\sum_i\E[v_i^1g_i^2]=\D_{h_0}(f)$. Next average
\eqref{eq:terminal-pair-bound} over the outside spins and sum over
ordered pairs $i\ne j$. The signed and absolute pair costs become the
two quadratic forms in \eqref{eq:terminal-curvature}. The quartic cost
is bounded by
\[
 2C_6\,r_4(U)\sum_i\E[v_i^1g_i^2].
\]
Increasing $C$ proves the lemma.
\end{proof}

The curvature inequality is useful when the few coordinates with appreciable
conditional variance occupy a small principal submatrix. Bandeira, El Alaoui,
and R\"odder~\cite[Theorem~1.1]{BER2026} combine uniform bounds on small
principal submatrices with a sufficiently large constant field to prove
polynomial mixing. Here the matrix bounds enter the signed curvature form
directly. We next obtain matrix control uniform over every such coordinate
set; the set will later depend on the spin configuration.

\begin{lemma}[Small principal submatrices]
\label{lem:terminal-matrices}
Fix $j\in(0,1)$.  There is a finite constant $M$, depending only on $j$,
with the following property.  For every $\varepsilon>0$ there is
$\alpha\in(0,1/2)$ such that, with probability tending to one over $J$,
\begin{gather}
 \max_{i,j}|J_{ij}|\le n^{-1/4},\qquad
 \|J\|,\|U\|,\|Q\|\le M,\qquad
 r_4(U)\le Mn^{-1/2},
 \label{eq:terminal-global-matrix}\displaybreak[1]\\
 \|U_{SS}\|\le\varepsilon,\qquad
 \|Q_{SS}\|\le\varepsilon
 \quad\text{for every }S\subseteq\{1,\ldots,n\}
                       \text{ with }|S|\le\alpha n.
 \label{eq:terminal-sparse-matrix}
\end{gather}
The same $M$ works for all $\varepsilon$.
\end{lemma}

\begin{proof}
For any deterministic unit vector $z$, the centered Gaussian variable
$z^{\mathsf T}Jz$ has variance
\[
 \frac{4j}{n}\sum_{i<k}z_i^2z_k^2\le\frac{2j}{n}.
\]
Thus
\begin{equation}
 \Prob\{|z^{\mathsf T}Jz|>t\}\le
                 2\exp\{-nt^2/(4j)\}.
 \label{eq:terminal-quadratic-tail}
\end{equation}
To pass from fixed quadratic forms to norms, choose a $1/4$-net of the
unit sphere in $\R^k$ with at most $9^k$ points. Every symmetric
$k\times k$ matrix $B$ then satisfies
\[
 \|B\|\le2\max_{z\text{ in the net}}|z^{\mathsf T}Bz|;
\]
to see this, approximate a unit vector maximizing the absolute quadratic
form and bound the change by $2\|B\|/4$.  Taking $k=n$ in
\eqref{eq:terminal-quadratic-tail} proves $\|J\|\le M_0$ with probability
at least $1-e^{-c n}$ for a sufficiently large constant $M_0=M_0(j)$.

For a fixed row $i$ and a fixed set $S$ of at most $k$ indices, Gaussian
independence within that row gives
\begin{equation}
 \Prob\left\{\sum_{\ell\in S}J_{i\ell}^2>t\right\}
 \le \exp\left\{-\frac{nt}{4j}+\frac{k}{2}\log2\right\}.
 \label{eq:terminal-row-tail}
\end{equation}
Use exponential Markov with parameter $n/(4j)$ to obtain this estimate:
each nonzero Gaussian square contributes the factor $\sqrt2$ to the
moment generating function. Choosing the full index set and taking a
union bound over rows bounds every row square sum by a fixed
$M_1=M_1(j)$, with failure at most $e^{-cn}$. Separately, the Gaussian
entry tail and a union bound give
\[
 \Prob\{\max_{i,\ell}|J_{i\ell}|>n^{-1/4}\}
 \le 2n^2\exp\{-\sqrt n/(2j)\}=o(1).
\]
Since $|\tanh z|\le|z|$, these bounds imply $\|Q\|\le M_1$ and
$r_4(U)\le M_1n^{-1/2}$, using the maximal absolute row-sum bound for
the operator norm of a symmetric matrix.  Moreover,
\[
 |\tanh z-z|\le |z|^3/3,
 \qquad
 \|U-J\|\le \frac13\max_i\sum_\ell|J_{i\ell}|^3
                \le\frac{M_1}{3}n^{-1/4}.
\]
The scalar inequality follows by integrating
$|1-\sech^2 z|=\tanh^2 z\le z^2$.  This proves
\eqref{eq:terminal-global-matrix}, with a fixed $M$.

The global bounds have fixed $M$; we now select the small-set fraction
$\alpha$. Put $k=\lfloor\alpha n\rfloor$. Bounding sets of size exactly
$k$ suffices, because a smaller set can be enlarged and restriction cannot
increase a principal-submatrix norm. The $\binom nk$ possible sets have
logarithmic count at most $n\mathfrak h(k/n)$, where
\[
 \mathfrak h(t)=-t\log t-(1-t)\log(1-t),
\]
and $\mathfrak h(t)\to0$ as $t\downarrow0$.  Applying the sphere-net
argument on each support and using \eqref{eq:terminal-quadratic-tail}
gives
\[
 \Prob\left\{\max_{|S|=k}\|J_{SS}\|>\varepsilon/2\right\}
 \le 2\binom nk9^k
                  \exp\{-n\varepsilon^2/(64j)\}.
\]
Choose $\alpha>0$ sufficiently small that this tends to zero
exponentially in $n$.  Then $\|U-J\|=o(1)$ proves the assertion for $U$.

Finally, \eqref{eq:terminal-row-tail}, with $t=\varepsilon$, and a union
bound over all rows and all $k$-element sets give
\[
 \Prob\left\{\max_i\max_{|S|=k}\sum_{\ell\in S}J_{i\ell}^2
                                     >\varepsilon\right\}
 \le n\binom nk
       \exp\left\{-\frac{n\varepsilon}{4j}+\frac{k}{2}\log2\right\}.
\]
Decreasing $\alpha$ further makes this probability exponentially small.
On its complement, the maximal row sum of every $Q_{SS}$ is at most
$\varepsilon$, which proves the remaining assertion.
\end{proof}

We apply these matrix bounds to a noisy observation. The following statement
allows any initial spin law, with the interaction still drawn at a fixed
$j\in(0,1)$; the exceptional probability concerns the observation at
that fixed matrix.

\begin{proposition}[Terminal observation gap]
\label{prop:terminal-gap}
Fix $j\in(0,1)$.  There are finite $T>0$, $c>0$, and events
$\mathcal E_n^{\mathrm{term}}$ depending only on $J$, with
$\Prob_J(\mathcal E_n^{\mathrm{term}})\to1$, such that the following holds
for every sufficiently large $n$ and every $J\in\mathcal E_n^{\mathrm{term}}$.
Let $X$ have any probability law on $\{-1,1\}^n$, possibly depending on
$J$, and let $Z$ be an independent standard Gaussian vector.  Set
$Y=TX+\sqrt T\,Z$.  Then, conditionally on $J$, except on an event of
probability at most $e^{-cn}$, the inequality
\begin{equation}
 \Var_{\mu_Y}(f)\le 2\D_Y(f)
 \label{eq:terminal-poincare}
\end{equation}
holds simultaneously for every real function $f$ on the cube.
All constants and matrix events are independent of the law of $X$.
\end{proposition}

\begin{proof}
Choose the tolerances in the order needed to absorb the curvature costs.
Let $C$ be the constant in Lemma~\ref{lem:curvature}, and choose
$\varepsilon>0$ with $(1+C)\varepsilon\le1/8$.
Lemma~\ref{lem:terminal-matrices} supplies $\alpha$ and the fixed global
bound $M$. Finally choose the fixed observation time $T$ large enough that
\begin{equation}
 \frac{16M^2}{T^2}\le\frac\alpha2,
 \qquad
 (1+C)M\bigl(2\sech(T/4)+\sech^2(T/4)\bigr)\le\frac18,
 \qquad
 e^{-T/8}\le\frac{\alpha}{4(e-1)}.
 \label{eq:terminal-T-choice}
\end{equation}
Let $\mathcal E_n^{\mathrm{term}}$ be the events supplied by that lemma.
For all sufficiently large $n$, their bounds also ensure
$\max|U_{ij}|\le u_0$ and $Cr_4(U)\le1/8$.

First count coordinates whose observation field remains small. Given
$X$, the indicators $\one_{\{|Y_i|<T/2\}}$ are independent; symmetry of
the Gaussian noise gives the same success probability $p_T$ for either
sign of $X_i$. The Gaussian tail bound yields
\[
 p_T=\Prob\{|T+\sqrt T Z_1|<T/2\}
       \le\Prob\{Z_1<-\sqrt T/2\}\le e^{-T/8}.
\]
For $B_Y=\#\{i:|Y_i|<T/2\}$, exponential Markov inequality yields
\[
 \Prob\{B_Y>\alpha n/2\mid X,J\}
 \le e^{-\alpha n/2}(1+p_T(e-1))^n
 \le e^{-\alpha n/4}.
\]
The same bound holds after averaging over $X$.

Suppose now that $B_Y\le\alpha n/2$.  For every spin configuration $x$,
$\|Jx\|^2\le M^2n$, and consequently
\[
 \#\{i:|(Jx)_i|>T/4\}\le\frac{16M^2n}{T^2}\le\frac\alpha2 n.
\]
It follows, simultaneously for all $x$, that
\begin{equation}
 S(x)=\{i:|Y_i+(Jx)_i|<T/4\}
 \quad\text{satisfies}\quad |S(x)|\le\alpha n.
 \label{eq:terminal-small-field-set}
\end{equation}
At this fixed field $Y$, set
$q_f(x)=\sum_i v_i^1(x)g_i(x)^2$ and
$\eta=\sech^2(T/4)$.  Since $0<v_i^1\le1$, one has pointwise
\[
 \|a\|^2\le q_f(x),\qquad
 \|a_{S(x)^c}\|^2\le\eta q_f(x).
\]
We may now use $S(x)$ even though it depends on the configuration,
because \eqref{eq:terminal-sparse-matrix} holds for every small set on
one matrix event. Split the quadratic form between $S(x)$ and its
complement. The small-set bound handles its first block, and the global
norm handles the cross and complementary terms, giving
\[
 a^{\mathsf T}Ua
 \le\bigl[\varepsilon+M(2\sqrt\eta+\eta)\bigr]q_f(x).
\]
The identical argument, with $|a|$ and $Q$, gives the same upper bound
for $|a|^{\mathsf T}Q|a|$.  Lemma~\ref{lem:curvature} and
\eqref{eq:terminal-T-choice} therefore imply
\begin{align*}
 \E_{\mu_Y}[(Hf)^2]
 &\ge\left[1-Cr_4(U)
       -(1+C)\{\varepsilon+M(2\sqrt\eta+\eta)\}\right]\D_Y(f)\\
 &\ge\frac12\D_Y(f).
\end{align*}

It remains to translate curvature into a gap. The kernel of the
self-adjoint operator $H$ contains exactly the constants: vanishing
$\E[fHf]$ forces every one-spin conditional variance to vanish, and the
strict positivity of the Gibbs law propagates equality across every
cube edge. For an eigenfunction with positive eigenvalue $\lambda$,
the curvature estimate and \eqref{eq:terminal-generator-form} imply
$\lambda^2\ge\lambda/2$, hence $\lambda\ge1/2$. Expand in an
orthonormal eigenbasis to obtain \eqref{eq:terminal-poincare}.
The observation failure bound was $e^{-\alpha n/4}$, uniformly in the
initial spin law, so we may set $c=\alpha/4$.
\end{proof}
